\documentclass[11pt]{article}
\usepackage{amsmath,amssymb}
\usepackage[margin=1in]{geometry}
\usepackage{hyperref}
\emergencystretch=3em

\title{The stochastic Taylor tree in every positive dimension: coefficient continuity and ordered remainders (Programme S, tranche A1--A2)}
\author{Claude, with Daniel Murfet}
\date{2026-09-08}

\begin{document}
\maketitle
\sloppy

\begin{abstract}
Units 270--277 lift the stochastic clause of \texttt{thm:strataempiricalexpansion} from the
single-chart $d=2$ statements of the normal-block programme to arbitrary normal dimension, on top of
the completed Taylor tree. On the Banach space $E_b\times E_b$ of weighted-$\ell^1$ phase and amplitude
data ($\sum_\gamma|c_\gamma|b^{|\gamma|}<\infty$, modelled as $\ell^1$ over a disjoint union of
multi-index sets in rescaled coordinates), every canonical coefficient $C_{\mu,j}$ of
$N^{-\mu}(\log N)^j$ in the expansion of the original box integral is Lipschitz on balls, hence
continuous and Borel measurable (Headline~XXXIV, a weighted-$\ell^1$ substitute for the paper's
cited-but-unstated \texttt{prop:convergence}). Ordering the terms by $(\nu,q)\prec(\mu,j)$ iff
$\nu<\mu$ or ($\nu=\mu$, $q>j$), the ordered normalised remainders converge to $C_{\mu,j}$ uniformly on
data balls, and for measurable random data $X_n\Rightarrow Z$ in $E_b\times E_b$ and sample sizes
$N_n\to\infty$, finite coefficient vectors and ordered normalised remainders converge in distribution,
jointly for finitely many targets (Headline~XXXV). Reviews v27 and v28: pass. Eight units against a cap
of eighteen. No \texttt{sorry}, no additional \texttt{axiom}; 3405 jobs; pin \texttt{caa3727}.
\end{abstract}

\section{The design (Astra \#33)}
After the Taylor tree was frozen, the instruction to continue formalising the paper pointed at
\S4.3. Astra~\#33 ranked the stochastic Taylor tree first: (A1) continuity of the coefficient
functionals, (A2) convergence in distribution of coefficients and ordered remainders, with a hard
review after the third unit; then (A3) tangential integration and chart assembly and (A4) posterior
division behind statement gates. Two design points shaped everything: use a genuine Banach space with
the frozen \texttt{CoeffFamily} theorems as its analytic interface, and treat the constant phase
$a=\xi(0)$ as a perturbation in its own right --- the Stage~4 stability gate had held it fixed.

\section{A1: the data space and the Lipschitz estimate (units 270--272)}
\paragraph{Data space (u270).} \texttt{DataSpace d} is real $\ell^1$ over
$(\mathrm{Fin}\,d\to\mathbb N)\sqcup(\mathrm{Fin}\,d\to\mathbb N)$. Its raw coordinates are exactly the
\emph{rescaled unit-box families} $\mathrm{scale}\,c_\xi\,b$, $\mathrm{scale}\,c_\eta\,b$ of the Taylor
tree, so the box families are recovered by dividing by $b^{|\gamma|}$ and the norm is the sum of the two
masses. The canonical map \texttt{dataBoxCoeff} is \texttt{boxCoeff} of the box families; through
\texttt{boxSpectralSum\_eq} it is literally the coefficient of $N^{-\mu}(\log N)^j$ for the original
integral, including the binomial conversion of $\log(Nb^{2|k|})$.
\paragraph{The gate (u271).} With $f=c_\eta*J^{*p}$, $J$ the constant-free phase, the coefficient series
term is $\beta^p/p!\;T_p(a;f)$. The difference of two data splits as
$[T_p(a';f')-T_p(a;f')]+T_p(a;f'-f)$. The first piece is new: by the mean value theorem on $[-R,R]$ and
$\partial_a\,\mathrm{fluctMoment}=\beta\,\mathrm{fluctMoment}(p{+}1)$,
$|\mathrm{fluctMoment}(a)-\mathrm{fluctMoment}(a')|\le\beta M_{\mu,n,p+1}(R)|a-a'|$, which propagates through
the kernel $S_p$ and the functional $T_p$. The second piece uses linearity and the mass algebra of the
Cauchy product, $\operatorname{mass}(f'-f)\le\Delta_\eta R^p+R\,pR^{p-1}\Delta_\xi$. The three series in
$p$ are summed in closed form by the Tonelli identity $\sum_p(\beta R)^p/p!\,M_{\nu,n,p}(R)=M_{\nu,n,0}(2R)$
and its shifted forms, giving the ballwise constant
$K_kD\bigl(2\beta RM_{\mu+1/2,n,0}(2R)+M_{\mu,n,0}(2R)\bigr)$.
\paragraph{Headline XXXIV (u272).} Transport to the data space (masses $\le\|x\|$, differences
$\le\|y-x\|$) and the finite log conversion give Lipschitz continuity on every closed ball, continuity
(the closed ball of radius $\|x\|+1$ is a neighbourhood), Borel measurability, and continuity of finite
coefficient vectors.

\section{A2: the deterministic core and the probability (units 273--277)}
\paragraph{Uniform constants (u273--274).} The cutoff constant is monotone in $|\xi(0)|$ and the two
masses, so on a data ball it is bounded by $\mathrm{cutoffBound}(R,R,R)$; the box integral itself is
Lipschitz in the data for fixed $N$ (the integrand is continuous on the closed cube by the Weierstrass
M-test, integrable on the box, and $|e^{A+s\xi}-e^{A+s\xi'}|\le e^{\beta\sqrt TR}\beta\sqrt T|\xi-\xi'|$
with $T=Nb^{2|k|}$), hence measurable.
\paragraph{Ordered remainders (u275).} With the cutoff $L=\mu+1$ the retained index set is the finite
$\Lambda_{\mu+1}\times\{0,\dots,n\}$; the predecessors of $(\mu,j)$ are its $\prec$-filter, independent
of the cutoff (u277). $Z-\sum_\prec-C_{\mu,j}N^{-\mu}(\log N)^j$ is the cutoff error plus the retained
terms with $\nu>\mu$ or ($\nu=\mu$, $q<j$); after normalisation these are
$O(N^{-(\nu-\mu)}(1+\log N)^n)$ and $O(1/\log N)$ with constants uniform on the ball, and the cutoff error
is $O(N^{-1}(1+\log N)^n)$. The explicit majorant is exported (\texttt{abs\_orderedRemainder\_sub\_le})
and gives \texttt{TendstoUniformlyOn} on the closed ball.
\paragraph{Headline XXXV (u276--277).} Coefficient vectors converge in distribution by Mathlib's
continuous mapping theorem. For the remainders, the norms $\|X_n\|$ are bounded in probability
(portmanteau on the closed sets $\{\|\cdot\|\ge m\}$, a generic copy of the $d=2$ argument), so uniform
convergence on balls gives convergence in probability of $R_{N_n}(X_n)-C(X_n)$ to zero, and
\texttt{tendstoInDistribution\_of\_tendstoInMeasure\_sub} (Slutsky) finishes. The joint version for
finitely many targets uses the sup norm on $\mathrm{Fin}\,m\to\mathbb R$ and
\texttt{Filter.eventually\_all}.

\section{Reviews and non-claims}
Review v27 (u270--272): pass, pass, qualified pass; it asked for the canonical cutoff to be restated
with \texttt{dataBoxCoeff}, for measurability of the integral and remainders, for the finite predecessor
set and descending-log ordering, for uniform convergence on balls with thresholds, and for the
boundedness-in-probability and Slutsky steps --- all done in u273--277. Review v28 (u273--276): pass on
all four units, ``no mathematical blocker''. Non-claims: the result is conditional on convergence in
distribution of the data in the weighted-$\ell^1$ topology at the fixed box radius $b$ (nothing is
derived from Hypothesis~I; no common analytic radius is manufactured for random phases); one chart, no
tangential integration or assembly; no Gaussianity of the limit; deterministic sample sizes
$N_n\ge0$ in the paper's convention; the normalised quotient is meaningful above the thresholds
$N\ge e$, $Nb^{2|k|}\ge1$.

\section{Lessons}
\begin{itemize}
\item \textbf{Use the induction structure you already have.} Every estimate here is a transport of a
frozen Taylor-tree theorem through the identification ``raw $\ell^1$ coordinates $=$ rescaled unit-box
families''; no new analysis of the integral was needed except the mean-value step in $\xi(0)$.
\item \textbf{Mean value beats differentiation under the integral.} The phase perturbation is a
one-line consequence of the existing derivative theorem for \texttt{fluctMoment}.
\item \textbf{Deterministic first, probability last.} Stating the uniform-on-balls theorem separately
made the probability step a two-lemma composition and let the reviewer check the analysis on its own.
\item \textbf{Tooling.} In this repository \texttt{lake env lean} cannot see the artifact-cached
package oleans; build single modules with \texttt{lake build Grammar.X}. Long-line lint must gate the
commit (three fix-up commits were needed before the gate was made mechanical). GPT consults must run in
the foreground on this host.
\item \textbf{Mathlib's probability API is enough.} \texttt{TendstoInDistribution} with its continuous
mapping, Slutsky and portmanteau lemmas covered A2 without new infrastructure.
\end{itemize}
\end{document}
